\section{Conclusion}\label{sec:conclusion}
%-------------------------------------------------------------------------------


% Comparison with standard strategies
The \rnn{} has an innovative design for electron selection based on
calorimetry information. With the
goal of satisfying specific trigger system demands, feature extraction
from concentric ring energy sums provides an efficient description, which is quickly derived through simple operations. 
Neural-network models are employed to exploit non-linear correlations between rings, and the optimization of this strategy 
culminated in the deployment of a set of models specific to each region of pseudorapidity and transverse energy. 


In the second half of 2017, the \rnn{} algorithm replaced a cut-based strategy
in the \fastcalo{} decision step of the \hlt{} and became the baseline algorithm
for triggering events containing electrons with \ET above \SI{15}{\GeV}, 
as part of the ATLAS online system's CPU-time reduction campaign. 
Furthermore, it was the first time a neural-network method was employed as a 
baseline event-selection algorithm in the ATLAS trigger system.


During 2017 operation, with the \rnn{} algorithm trained on simulated data, 
fake-electron trigger rates were reduced by a factor of 13 in the \fastcalo{} step 
and by a factor of 2 overall in the \hlt{} with respect to the offline likelihood algorithm, while
achieving high electron-selection efficiencies very similar to those of relevant single-electron triggers without the \rnn{}. 
For 2018, the \rnn{} was trained with collision
data, resulting in an estimated additional reduction factor of 1.19--1.69 for fake-electron triggers in the \fastcalo{} step with respect to 2017. 
The \rnn{} was also studied from an offline likelihood perspective to obtain additional insights into its behaviour by using two equivalent single-electron triggers with and without the \rnn{}. The \enquote{quadrant analysis} showed a high level of agreement between the two triggers, with the \rnn{} trigger showing a slight bias towards the signal region in the disagreement profiles of some calorimetry-based variables. Finally, the \enquote{homogeneity tests} between the two groups of data resulted in similar $p$-values. In other words, the fluctuations in the profiles are mainly statistical, with no sign of systematic effects due to the trigger configuration.



Physics beyond the Standard Model is getting more attention at the LHC and is expected to be part of the main focus in the data-taking phases to
come. At the beginning of \RunThr, the \rnn was updated to include events in which electrons have smaller angular separations (such as in boosted topologies), particularly limiting the effect of other contributions at
the edge of the feature extraction window, which can result from these physics
processes. In addition, the \rnn was extended into the \ET region below 15~GeV, motivated by the high CPU demands of triggers in this kinematic region.





\begin{comment}



    


% Quadrant and agreement analysis
The \rnn{} was also studied through the offline likelihood perspective 
to bring additional insights on its behavior. By
taking advantage of the low dimensionality, low correlation and high
interpretability power of the variables employed in the likelihood, we derived
the profiles of the disjoint decision cases of the 2017 duplicate trigger pair.
The choice for the duplicated pair considered the typical scenario employed for
analysis with the setup that would maximize the disagreement between both
configurations: with and without the \rnn{}. The proposed quadrant analysis
allowed to observe high agreement between both triggers with slight shifts
towards the signal region in the disagreement profiles of some calorimetry-based
variables. By checking the agreement between the trigger pair for the derivation
of the likelihood pdfs, the alteration using the full 2017 data was much lower
than estimated statistical fluctuations. The residuals were found to be bounded
by \SI{0.2}{$\sigma$} for all variables. When considering only the
disagreement profiles, homogeneity hypothesis at a \SI{5}{\%} level was not
rejected for the calorimetry variables, except for \reta{} and \rhad{} despite
using the most populated regions. For these variables, a shift towards signal
region was also observed. Hence, the \rnn{} had negligible impact in the
sensitive variables employed for physics analysis and in the offline selection
operation.

Finally, the presented results show that it is possible to tighten the requirements for online operation while still resulting in high agreement with the offline
selection. More generally, it demonstrates that designing a solution from scratch using the application specificities can provide new effective
strategies. We hope that it can serve as motivation for other developments in
HEP, particularly when considering the trigger applications.

\section{Outlook}


Preliminary evaluations have shown the potential of the \rnn{} algorithm for
triggering on electrons with $\et<\SI{15}{\GeV}$. This is particularly motivating given that triggers in this kinematic region are very
CPU demanding. Additionally, they have high output rate and are usually
prescaled, thus \rnn{} might also allow to collect additional data for those
triggers. For the beginning of Run 3 (2022), \rnn{} has been extended to all triggers below $\et<\SI{15}{\GeV}$ where a reduction of at least half of the fake rate in \fastcalo was observed.


Photons are other physics objects that may be triggered using the Ringer approach. Triggers combining photons to other physics objects are quite CPU demanding, thus motivating studies to evaluate whether additional Ringer developments will contribute to increased performance.
When considering pure photon triggers, the ring sums with linear classifiers can be investigated in order to keep the analysis strategies employed in channels as $\text{H}\rightarrow Z\gamma$. Another possibility for those triggers is to
employ the \rnn{} for increasing the trigger efficiency. A possible
setup is to duplicate the photon triggers so that the \rnn{} recovers
interesting events.  In this hybrid menu configuration, the standard trigger
would be available for those analyses so that they may benefit from high
interpretability and customization, whereas the \rnn{} trigger might be employed
for those analyses demanding more statistics. It is expected that this
configuration will demand minor additional resources but may allow higher
statistical significance of the results of physics analysis.


% No-had
Another interesting possibility is to employ the \rnn{} accessing only
electromagnetic information for triggers resulting in high readout rate from
hadronic calorimeter cells. The particular setup may even consider employing a
two step decision on the \fastcalo{}, a preliminary one based solely on
electromagnetic information followed by the full \rnn{} decision.





% Calibration
The shower description by ring sums goes beyond capturing discriminant
information and may be used for calibration of the energy of electrons and
photons. While preliminary results with ring sums are motivating, we expect the aforementioned
improvements in the ring description to be exceptionally important for
calibration.

% Offline developments
The \rnn{} software has been extended to the offline framework. Currently,
the offline ring description is available for all electrons with $\et>\SI{14}{\GeV}$. This is an on-going activity.


\end{comment}
